\documentclass[12pt,a4paper]{report}

\usepackage{elearning_i18n}
\usepackage{microtype}

\ElearnPlatformName{Delta}
\ElearnOrganization{Delta Education}
\ElearnDocumentType{Specyfikacja funkcjonalna platformy e-learningowej}
\ElearnDocumentTitle{Analiza interfejsu i wymagania funkcjonalne}
\ElearnDocumentSubtitle{Wnioski z przeglądu makiet, komentarze produktowe i kierunki implementacji}
\ElearnDocumentVersion{1.0}
\ElearnDocumentStatus{Dokument roboczy po przeglądzie makiet}
\ElearnDocumentOwner{Zespół Delta}
\ElearnDocumentAudience{Właściciele produktu, projektanci UX/UI, programiści i autorzy treści}
\ElearnDocumentDate{22 lipca 2026}
\ElearnDocumentKeywords{Delta, e-learning, matematyka, wymagania funkcjonalne, UX, kursy online}
\ElearnConfidentiality{Dokument wewnętrzny}

% Logo path is relative to this main .tex file / compilation directory.
% Replace it with the actual Delta logo path.
\ElearnLogoPath{assets/logo_large.png}
% Width only, expressed as a percentage of \textwidth.
% The logo aspect ratio is preserved automatically.
\ElearnLogoWidthPercent{40}

\ElearnWatermarkText{DELTA}
\ElearnWatermarkColor{black!4}
\ElearnWatermarkAngle{35}

\begin{document}
\selectlanguage{polish}
\ElearnFrontMatter
\chapter{Cel i zakres dokumentu}

Dokument porządkuje materiały z pliku źródłowego zawierającego makiety interfejsu oraz komentarze recenzenckie. Zrzuty ekranu zostały zachowane jako materiał referencyjny, a komentarze umieszczono przy odpowiadających im obszarach funkcjonalnych. Na ich podstawie sformułowano wymagania, decyzje produktowe i kwestie wymagające doprecyzowania.

\section{Najważniejsze decyzje produktowe}
\begin{itemize}
  \item Konto użytkownika i potwierdzenie adresu e-mail są podstawą dostępu do większości funkcji platformy.
  \item Główna zakładka z materiałami ma nosić nazwę "Nauka", a jej struktura powinna obsługiwać kolejne przedmioty i poziomy edukacji.
  \item Zadania rozwiązane poprawnie są usuwane z puli, natomiast zadania rozwiązane błędnie wracają do późniejszej powtórki.
  \item Lekcje i rozwiązania zadań mają być prezentowane jako materiały z narracją głosową; notatki mają być dostępne jako PDF.
  \item Profil użytkownika ma zawierać informacje o koncie, aktywnym kursie, terminie ważności i kodach rabatowych.
  \item Moduł maturalny ma koncentrować się na zadaniach, w tym zadaniach z arkuszy, bez obietnic "pewniaków maturalnych".
  \item System autorski musi umożliwiać przygotowywanie rysunków geometrycznych oraz wykresów funkcji.
\end{itemize}
\section{Materiały źródłowe}

Dokument źródłowy składał się z 19 osadzonych obrazów i 17 komentarzy. Jeden komentarz był pusty. Obrazy są prezentowane w kolejności logicznej odpowiadającej przepływowi użytkownika oraz strukturze produktu.

\chapter{Rejestracja, konto i początkowa konfiguracja}
\section{Wybór dni nauki i dostęp do platformy}
\ElearnScreenshot[0.72]{assets/screen_01.png}{Ekran konfiguracji dni, w których użytkownik chce otrzymywać przypomnienia o nauce.}
\begin{ElearnObservedUI}
\begin{itemize}
  \item Adres e-mail użytkownika w nagłówku.
  \item Wybór dni tygodnia od poniedziałku do niedzieli.
  \item Informacja o częstotliwości nauki i komunikat motywacyjny.
  \item Przycisk "Kontynuuj".
\end{itemize}
\end{ElearnObservedUI}
\begin{ElearnReviewComment}[title={Komentarz nr 0 -- Anna Kogut, 2026-07-22}]
Ustalanie ile dni w tygodniu mamy wysłać przypomnienie o nauce

Przede wszystkim żeby cokolwiek zobaczyć (oprócz cen kursów i dostępu do zadań) trzeba utworzyć konto, podając imię i nazwisko , adres email i utworzyć hasło.

Potwierdzenie rejestracji przyjdzie na podany adres emailowy i trzeba będzie potwierdzić.

Wchodząc w zakładkę kursu i chcąc kupić kurs - trzeba podać adres emailowy. Jeżeli adres emailowy nie zostanie rozpoznany w bazie danych, przekieruje do utworzenia konta i potwierdzenia w emailu. Wtedy użytkownik może przejść ponownie na kurs i kupić odpwiedni pakiet.
\end{ElearnReviewComment}
\begin{ElearnRequirement}[title={Wymagania wynikające z komentarza}]
\begin{itemize}
  \item Użytkownik określa, w które dni tygodnia platforma ma wysyłać przypomnienia o nauce.
  \item Dostęp do większości treści i funkcji wymaga utworzenia konta z imieniem, nazwiskiem, adresem e-mail i hasłem.
  \item Rejestracja wymaga potwierdzenia adresu e-mail za pomocą wiadomości aktywacyjnej.
  \item Przy zakupie kursu system sprawdza adres e-mail. Nieznany adres uruchamia proces rejestracji i potwierdzenia konta.
  \item Po aktywacji konta użytkownik wraca do oferty i wybiera właściwy pakiet.
\end{itemize}
\end{ElearnRequirement}
\begin{ElearnWarning}
Do doprecyzowania pozostaje zakres publicznego dostępu bez konta. Komentarz źródłowy wskazuje ceny kursów i dostęp do zadań jako wyjątki, ale docelowe ograniczenia powinny zostać jednoznacznie opisane w regułach autoryzacji.
\end{ElearnWarning}
\clearpage

\chapter{Plan nauki i powiadomienia}
\section{Dzienny cel}
\ElearnScreenshot[0.92]{assets/screen_02.png}{Wariant ekranu ustawiania dziennego celu w widoku przeglądarkowym.}
\ElearnScreenshot[0.67]{assets/screen_04.png}{Wariant mobilny ustawiania dziennego celu na siedem lekcji i zadań.}
\begin{ElearnObservedUI}
\begin{itemize}
  \item Kontrolka minus / wartość / plus służąca do zmiany dziennego celu.
  \item Komunikat motywacyjny zależny od wybranej wartości.
  \item Przycisk zapisania ustawienia.
\end{itemize}
\end{ElearnObservedUI}
\begin{ElearnReviewComment}[title={Komentarz nr 1 -- Anna Kogut, 2026-07-22}]
Ustalanie ile tematów dziennie uczeń chce przerobić, ile zadań - tutaj tylko zadania a nie tematy. Jak uczeń chce przejść do innego tematu, jego wola. Musi zezwolić o przesłaniu powiadomienia na email

Komentarz odnosi się do ustawiania dziennego celu oraz zgody na powiadomienia.
\end{ElearnReviewComment}
\begin{ElearnRequirement}[title={Wymagania wynikające z komentarza}]
\begin{itemize}
  \item Dzienny cel powinien dotyczyć liczby zadań do wykonania, a nie liczby tematów obowiązkowo przechodzonych przez ucznia.
  \item Uczeń zachowuje swobodę przechodzenia do innych tematów niezależnie od ustawionego celu.
  \item Włączenie przypomnień e-mail wymaga świadomej zgody użytkownika.
\end{itemize}
\end{ElearnRequirement}
\section{Zgoda na powiadomienia}
\ElearnScreenshot[0.68]{assets/screen_03.png}{Prośba o zezwolenie na wysyłanie przypomnień o nauce.}
\begin{ElearnRequirement}[title={Uzupełniające wymagania UX}]
\begin{itemize}
  \item System powinien wyjaśniać cel powiadomień przed poproszeniem o zgodę.
  \item Odmowa nie może blokować korzystania z kursu; zgodę można ponownie zmienić w profilu lub ustawieniach.
\end{itemize}
\end{ElearnRequirement}
\section{Podsumowanie planu nauki}
\ElearnScreenshot[0.62]{assets/screen_05.png}{Podsumowanie osobistego planu nauczania: dzienny cel, dni nauki i godzina przypomnienia.}
\begin{ElearnReviewComment}[title={Komentarz nr 2 -- Anna Kogut, 2026-07-22}]
Podsumowanie

Komentarz pełni rolę oznaczenia ekranu podsumowującego i nie zawiera dodatkowych zaleceń.
\end{ElearnReviewComment}
\begin{ElearnRequirement}[title={Wymagania wynikające z komentarza}]
\begin{itemize}
  \item Podsumowanie pokazuje dzienny cel, dni nauki i godzinę przypomnienia.
  \item Każdy parametr powinien być edytowalny bez ponownego przechodzenia całego onboardingu.
  \item Po zatwierdzeniu użytkownik przechodzi bezpośrednio do kursu.
\end{itemize}
\end{ElearnRequirement}
\clearpage
\chapter{Strona główna, rabaty i funkcje dodatkowe}
\section{Panel główny i kody rabatowe}
\ElearnScreenshot[0.78]{assets/screen_06.png}{Panel główny z serią nauki, kodem rabatowym i odnośnikami społecznościowymi.}
\begin{ElearnObservedUI}
\begin{itemize}
  \item Licznik kolejnych dni nauki.
  \item Karta z kodem rabatowym i ograniczeniem czasowym.
  \item Odnośniki do grup społecznościowych.
\end{itemize}
\end{ElearnObservedUI}
\begin{ElearnReviewComment}[title={Komentarz nr 3 -- Anna Kogut, 2026-07-22}]
Po wejściu w zakladke kursy możemy podać kod do rabatu. Natomiast ograniczenie ze kod wygaśnie za 2 dni jest wymuszeniem kupienia kolejnego kursu.

Uczeń wykupuje miesięczny abonanment - dostanie 5\% rabatu od ceny wyjściowej (do jednorazowego wykorzystania, który będzie widniał w jego profilu (koło zębate).

Uczeń wykupuje 3 m-czny abonament dostanie 10\% rabatu.

Uczeń wykupi 6 m-czny kurs dostanie 15\%rabatu.

Uczeń wykupi roczny abonament - dostanie 20\% pod warunkiem, że wykupi w kolenym roku roczny abonament.
\end{ElearnReviewComment}
\begin{ElearnRequirement}[title={Wymagania wynikające z komentarza}]
\begin{itemize}
  \item Nie należy stosować krótkiego, dwudniowego terminu ważności kodu jako mechanizmu presji zakupowej.
  \item Zakup miesięcznego abonamentu daje jednorazowy rabat 5\% od ceny wyjściowej.
  \item Zakup abonamentu trzymiesięcznego daje rabat 10\%, sześciomiesięcznego 15\%, a rocznego 20\% przy odnowieniu rocznego abonamentu w kolejnym roku.
  \item Niewykorzystane kody rabatowe są widoczne w profilu użytkownika.
  \item Regulamin powinien definiować sposób naliczania, ważność i możliwość łączenia rabatów.
\end{itemize}
\end{ElearnRequirement}
\section{Opinie, kontakt i zapisane materiały}
\ElearnScreenshot[0.70]{assets/screen_07.png}{Karty funkcji dodatkowych: głosowanie na ulepszenia, obsługa klienta, zapisane materiały i aplikacja mobilna.}
\begin{ElearnReviewComment}[title={Komentarz nr 4 -- Anna Kogut, 2026-07-22}]
Zagłosuje na ulepszenia - super ...trzeba pomyśleć o ankiecie. Pytanie gdzie te dane będą gromadzone.

Czat z obsługą klienta zamieniłabym na kontakt - możliwość napisania emaila. Formularz z imieniem , adres emailowy i okienko do napisania tekstu w czym problem i wyślij

Zapisane - usunęłabym, bo notatki z lekcji będzie można wydrukować lub sciągnąć i zapisać sobie na własny komputer
\end{ElearnReviewComment}
\begin{ElearnRequirement}[title={Wymagania wynikające z komentarza}]
\begin{itemize}
  \item Funkcja "Zagłosuj na ulepszenie" powinna prowadzić do ankiety lub formularza opinii; trzeba określić miejsce przechowywania i analizowania danych.
  \item Czat z obsługą klienta należy zastąpić formularzem kontaktowym zawierającym imię, adres e-mail, treść problemu i przycisk wysłania.
  \item Oddzielna sekcja "Zapisane" nie jest wymagana, ponieważ notatki mają być pobierane lub drukowane i przechowywane lokalnie przez użytkownika.
  \item Odnośnik do aplikacji mobilnej może pozostać jako element przyszłej wersji produktu.
\end{itemize}
\end{ElearnRequirement}
\clearpage
\chapter{Struktura sekcji Nauka}
\section{Nawigacja po poziomach i działach}
\ElearnScreenshot[0.92]{assets/screen_08.png}{Widok hierarchii materiałów z podziałem na klasy i działy.}
\begin{ElearnReviewComment}[title={Komentarz nr 5 -- Anna Kogut, 2026-07-22}]
Działy nazwałabym "nauka" bo w przyszłości dodamy tam rózne sekcje (matematyka w szkole średniej, matematyka w szkole podstawowej, fizyka itd....)

Na początek z zakładce NAUKA umieścimy trzy podfoldery "matematyka - szkoła średnia" "matura - poziom podstawowy" oraz "matura - poziom rozszerzony"

W szkole średniej będzie jak tutaj podział na klasy i po rozwinięciu pokażą się rozdziały.

Klikając na poszczególny rozdział będą widoczne zadagnienia do omówienia oraz zadania do rozwiązania ---> patrz niżej
\end{ElearnReviewComment}
\begin{ElearnRequirement}[title={Wymagania wynikające z komentarza}]
\begin{itemize}
  \item Zakładkę "Działy" należy nazwać "Nauka", aby mogła obejmować wiele przedmiotów i poziomów edukacji.
  \item Pierwsza wersja sekcji Nauka zawiera trzy obszary: "Matematyka - szkoła średnia", "Matura - poziom podstawowy" oraz "Matura - poziom rozszerzony".
  \item W obszarze szkoły średniej materiały są podzielone na klasy, następnie na działy i zagadnienia.
  \item Po wejściu do działu użytkownik widzi zagadnienia do omówienia i zadania do rozwiązania.
  \item Architektura nawigacji powinna być rozszerzalna o kolejne przedmioty, np. fizykę.
\end{itemize}
\end{ElearnRequirement}
\section{Widok zagadnienia i pula zadań}
\ElearnScreenshot[0.92]{assets/screen_09.png}{Widok zagadnienia z notatkami, lekcją wprowadzającą, zadaniami i generowaniem sprawdzianu.}
\begin{ElearnReviewComment}[title={Komentarz nr 6 -- Anna Kogut, 2026-07-22}]
Po wejściu w dany temat - zaganienie, pojawi się Lekcja oraz zadania. Na końcu dałabym notatki do wydrukowania

Przy zadanaich też umieściłabym liczbę zadań do rozwiązania (liczba wszystkich zadań do rozwiązania) z wyświeleniem ile zostało zrobionych poprawnie. Błędne wracają do puli zadań do rozwiązania. Nie wyświetlamy błędnych. Zrobione dobrze nie wracają do puli. Uczeń mając poprawnie rozwiazane 5 zadań z 135, będzie miał 130 do rozwiązania.

Wygeneruj sprawdzian - na razie to usuwamy.

Generowanie zestawu zadań zrobiłabym w "matura" gdzie uczeń będzie mógł się sprawdzić z całego materiału ale o tym później
\end{ElearnReviewComment}
\begin{ElearnRequirement}[title={Wymagania wynikające z komentarza}]
\begin{itemize}
  \item Po wejściu w zagadnienie widoczne są: lekcja, zadania oraz na końcu notatki do pobrania lub wydruku.
  \item Sekcja zadań pokazuje całkowitą liczbę zadań oraz liczbę zadań rozwiązanych poprawnie.
  \item Zadanie rozwiązane błędnie wraca do puli, natomiast zadanie rozwiązane poprawnie nie pojawia się ponownie w bieżącej puli.
  \item Licznik "do rozwiązania" jest obliczany jako liczba wszystkich zadań pomniejszona o liczbę poprawnie rozwiązanych.
  \item Funkcję "Wygeneruj sprawdzian" należy usunąć z pierwszej wersji.
  \item Generowanie zestawów zadań zostanie rozważone później w module maturalnym.
\end{itemize}
\end{ElearnRequirement}
\clearpage
\chapter{Materiały edukacyjne}
\section{Notatki do lekcji}
\ElearnScreenshot[0.92]{assets/screen_10.png}{Przykład notatki z lekcji prezentowanej w platformie.}
\begin{ElearnReviewComment}[title={Komentarz nr 7 -- Anna Kogut, 2026-07-22}]
Notatki będą w formie PDF do sciągnięcia lub wydrukowania - najważniejsze notatki z lekcji
\end{ElearnReviewComment}
\begin{ElearnRequirement}[title={Wymagania wynikające z komentarza}]
\begin{itemize}
  \item Notatki zawierają najważniejsze informacje z lekcji.
  \item Notatki są dostarczane jako pliki PDF.
  \item Użytkownik może pobrać plik lub wydrukować go bezpośrednio z platformy.
\end{itemize}
\end{ElearnRequirement}
\section{Wyjaśnienie rozwiązania zadania}
\ElearnScreenshot[0.70]{assets/screen_11.png}{Przykład materiału pokazującego ręczne rozwiązanie zadania.}
\begin{ElearnReviewComment}[title={Komentarz nr 8 -- Anna Kogut, 2026-07-22}]
Dane zadanie zawsze pod spodem odpowiedzi a b c d będzie miało plik w formie ??? z wyjaśnieniem jak to zadanie rozwiązać.. Tu dałambym plik o którym rozmawialiśmy, z głosem i ewentualnie czerwony kursor wskazujący w którym momencie opowiadam o zagadnieniu
\end{ElearnReviewComment}
\begin{ElearnRequirement}[title={Wymagania wynikające z komentarza}]
\begin{itemize}
  \item Każde zadanie powinno mieć powiązany materiał wyjaśniający sposób rozwiązania.
  \item Materiał powinien zawierać narrację głosową i wizualne wskazywanie omawianego fragmentu, np. czerwonym kursorem.
  \item Docelowy format pliku należy wybrać na etapie projektu technicznego; praktycznym kandydatem jest wideo MP4/WebM lub nagranie ekranu z warstwą audio.
  \item Odtwarzacz powinien działać w przeglądarce i na urządzeniach mobilnych.
\end{itemize}
\end{ElearnRequirement}
\clearpage
\chapter{Mechanika rozwiązywania zadań}
\section{Zadania zamknięte i informacja zwrotna}
\ElearnScreenshot[0.92]{assets/screen_12.png}{Zadanie wielokrotnego wyboru z zaznaczoną poprawną odpowiedzią i materiałem wideo pod zadaniem.}
\begin{ElearnReviewComment}[title={Komentarz nr 9 -- Anna Kogut, 2026-07-22}]
Każde zadanie ma odpowiedź a, b c d. Po kliknięciu w prawidłową odpowiedź, będzie się wyświetlało jak obok i dodałabym cos w rodzaju uśmiechu że uczeń robi postępy.

Gdy jednak zaznaczy źle, zadanie wróci do puli zadań do rozwiązania z tekstem "spróbuj ponownie za chwilę" nie koniecznie od razu

Wyświetli mu się wtedy plik w rozwiązaniem zadania
\end{ElearnReviewComment}
\begin{ElearnRequirement}[title={Wymagania wynikające z komentarza}]
\begin{itemize}
  \item Standardowe zadanie zamknięte zawiera odpowiedzi A, B, C i D.
  \item Po wskazaniu poprawnej odpowiedzi interfejs potwierdza wynik i prezentuje pozytywną, motywującą informację zwrotną.
  \item Po odpowiedzi błędnej zadanie wraca do puli i nie musi być pokazane ponownie natychmiast; komunikat może brzmieć "Spróbuj ponownie za chwilę".
  \item Po błędnej odpowiedzi użytkownik otrzymuje dostęp do materiału z rozwiązaniem i wyjaśnieniem.
  \item Stan zadania powinien być zapisany w historii użytkownika, ale widok bieżącej puli nie pokazuje osobnego licznika błędnych odpowiedzi.
\end{itemize}
\end{ElearnRequirement}
\clearpage
\chapter{Postępy użytkownika}
\section{Postęp według działów}
\ElearnScreenshot[0.92]{assets/screen_13.png}{Referencyjny widok historii aktywności i postępów użytkownika.}
\begin{ElearnReviewComment}[title={Komentarz nr 10 -- Anna Kogut, 2026-07-22}]
Wchodząc w zakładkę postępy widziałabym to tak:

Dane zaganienie ( klasa 1 - zbiory liczb)  i do niego przypisane ilość poprawnie rozwiązanych na ilość wszystkich zadań, bez dat i godzin

Bez przypisywania co było robione dziś czy wczoraj. Aby zaliczyć dział trzeba rozwiązać wszystkie zadania z tego działu poprawnie
\end{ElearnReviewComment}
\begin{ElearnRequirement}[title={Wymagania wynikające z komentarza}]
\begin{itemize}
  \item Widok postępów ma być agregowany według zagadnienia lub działu, np. "Klasa 1 - zbiory liczb".
  \item Dla każdego elementu należy pokazywać liczbę poprawnie rozwiązanych zadań w relacji do liczby wszystkich zadań.
  \item Nie jest wymagany dziennik aktywności z datami i godzinami ani podział na "dzisiaj" i "wczoraj".
  \item Dział zostaje zaliczony po poprawnym rozwiązaniu wszystkich przypisanych do niego zadań.
  \item Interfejs powinien jasno odróżniać postęp częściowy od pełnego zaliczenia.
\end{itemize}
\end{ElearnRequirement}
\clearpage
\chapter{Lekcje i prezentacja treści}
\section{Lista lekcji w dziale}
\ElearnScreenshot[0.92]{assets/screen_14.png}{Lista lekcji w dziale "Klasa 1: Wyrażenia algebraiczne".}
\begin{ElearnReviewComment}[title={Komentarz nr 11 -- Anna Kogut, 2026-07-22}]
Wchodzą w zakładkę Nauka ---> Matematyka - szkoła średnia ---> klasa 1 ---> dział "wyrażenia algebraiczne" ...widzę lekcje dotyczące tego działu

Po kliknięciu tematu mam to co na stronie 7 tego pliku , czyli lekcja wprowadzająca, pod spodem jest "zadania" i notatki do wydruku lub sciągnięcia
\end{ElearnReviewComment}
\begin{ElearnRequirement}[title={Wymagania wynikające z komentarza}]
\begin{itemize}
  \item Ścieżka nawigacji ma postać: Nauka -> Matematyka - szkoła średnia -> Klasa -> Dział -> Lekcja.
  \item Po wybraniu działu użytkownik widzi listę lekcji dotyczących tego działu.
  \item Po wejściu w temat strona zawiera lekcję wprowadzającą, zadania oraz notatki do pobrania lub wydruku.
  \item Nawigacja powinna pokazywać kontekst bieżącego działu i umożliwiać powrót o jeden poziom.
\end{itemize}
\end{ElearnRequirement}
\section{Format lekcji}
\ElearnScreenshot[0.92]{assets/screen_15.png}{Przykład lekcji w formie nagrania omawiającego notatkę i wykres funkcji.}
\begin{ElearnReviewComment}[title={Komentarz nr 12 -- Anna Kogut, 2026-07-22}]
Każda lekcja wolałabym żeby była w pliku jako nagranie z głosem z którym wyjaśniam dane zagadnienie
\end{ElearnReviewComment}
\begin{ElearnRequirement}[title={Wymagania wynikające z komentarza}]
\begin{itemize}
  \item Każda lekcja powinna być dostępna jako nagranie z narracją głosową nauczyciela.
  \item Nagranie może prezentować notatkę, rysunek, zapis rozumowania lub wykres.
  \item Odtwarzacz powinien zapewniać pauzę, przewijanie, zmianę głośności oraz działanie pełnoekranowe.
  \item W przyszłości warto przewidzieć napisy lub transkrypcję dla dostępności.
\end{itemize}
\end{ElearnRequirement}
\clearpage
\chapter{Profil użytkownika}
\section{Zakres danych w profilu}
\ElearnScreenshot[0.92]{assets/screen_16.png}{Referencyjny widok profilu i ustawień użytkownika.}
\begin{ElearnReviewComment}[title={Komentarz nr 13 -- Anna Kogut, 2026-07-22}]
Zakładka koło zębate to profil użtkownika w którym będzie:

- konto z jego adres email, resetuj hasło

- plan nauki bym nie wstawiała (pomijamy to)

- rodzaj kursu z którego obecnie korzysta z data ważności.

- twoje kody rabatowe
\end{ElearnReviewComment}
\begin{ElearnRequirement}[title={Wymagania wynikające z komentarza}]
\begin{itemize}
  \item Profil zawiera adres e-mail użytkownika oraz funkcję resetowania hasła.
  \item Sekcja "Plan nauki" nie jest wymagana w profilu i może zostać pominięta.
  \item Profil pokazuje rodzaj aktualnie wykupionego kursu lub abonamentu oraz datę jego ważności.
  \item Profil zawiera listę dostępnych kodów rabatowych i ich warunki wykorzystania.
  \item Należy zachować funkcję wylogowania oraz możliwość zarządzania zgodami na powiadomienia.
\end{itemize}
\end{ElearnRequirement}
\clearpage
\chapter{Narzędzia do przygotowywania zadań}
\section{Rysunki geometryczne i wykresy funkcji}
\ElearnScreenshot[0.73]{assets/screen_17.png}{Przykład zadania wymagającego rysunku bryły geometrycznej.}
\ElearnScreenshot[0.92]{assets/screen_18.png}{Przykład zewnętrznego edytora wykresów funkcji.}
\begin{ElearnReviewComment}[title={Komentarz nr 14 -- Anna Kogut, 2026-07-22}]
Niektóre zadania będą wymagały rysunków, więc potrzebny będzie edytor to tworzenia takich rysunków do zadań, albo np. Będę musiała narysować funcję do zadania, jakiś wykres

Na matemks.pl jest taki program do rysowania funkcji.
\end{ElearnReviewComment}
\begin{ElearnReviewComment}[title={Komentarz nr 15 -- Anna Kogut, 2026-07-22}]
Komentarz w dokumencie źródłowym nie zawierał treści.

Pusty komentarz występował obok materiału dotyczącego narzędzia do rysowania wykresów.
\end{ElearnReviewComment}
\begin{ElearnRequirement}[title={Wymagania wynikające z komentarza}]
\begin{itemize}
  \item Panel autorski musi umożliwiać dodawanie rysunków do treści zadań.
  \item Potrzebne jest narzędzie do przygotowywania prostych diagramów geometrycznych oraz wykresów funkcji.
  \item Możliwe są dwa warianty: wbudowany edytor albo integracja z zewnętrznym narzędziem i import gotowego obrazu.
  \item Obraz powinien mieć tekst alternatywny i zachowywać czytelność na ekranach mobilnych.
  \item Wymagania edytora należy opisać osobno: obsługiwane obiekty, osie, skala, etykiety, eksport SVG/PNG i przechowywanie plików źródłowych.
\end{itemize}
\end{ElearnRequirement}
\clearpage
\chapter{Moduł maturalny}
\section{Matura - poziom podstawowy i rozszerzony}
\ElearnScreenshot[0.92]{assets/screen_19.png}{Widok oferty kursu maturalnego z listą działów, liczbą lekcji i zadań.}
\begin{ElearnReviewComment}[title={Komentarz nr 16 -- Anna Kogut, 2026-07-22}]
W zakładce nauka będzie segment matura - podstawowa. Po wejściu w ten segment będą działy z podziałem na działy jak w działach we wszystkich klasach.

Po wejściu w dział będą zadania z danego działu (więc tu będzie mniej pracy , bo połączymy wszystkie zadania ze wszystkich lekcji z danego działu) oraz zadania dodatkowe z matur z arkuszy maturalnych

Zadania otwarte będą pomieszane z zamkniętymi.

Zadnia otwarte będą miały jedną odpowiedź, którą uczeń będzie musiał wpisać w okienko. Błędnie wpisana odpowiedź spowoduje, że zadanie wróci do puli zadań do wykonania.

W tym segmencie nie będzie lekcji. Tylko w przypadku błędnej odpowiedzi pojawi się rozwiązanie do tego zadania z wyjaśnieniem. Tutaj tez mamy mniej pracy, bo rozwiązania zadania z lekcji powielą się. Jedynie będę musiała zrobić rozwiżąnia z wyjaśnieniem do zadań z arkuszy.

U nas nie ma i będzie pewniaków maturalnych. Bo to jest bzdura. To tak jakby uczeń zagrał w toto lotka. Matury teraz są na zasadzie zadań, z każdego praktycznie działu z podstawy programowej, choćby jedno zadanie ale jest.

W trakcie rozwoju strony będę dodawać swoje indywidualne zadania.
\end{ElearnReviewComment}
\begin{ElearnRequirement}[title={Wymagania wynikające z komentarza}]
\begin{itemize}
  \item W sekcji Nauka dostępny jest segment "Matura - poziom podstawowy"; analogiczna struktura może zostać przygotowana dla poziomu rozszerzonego.
  \item Segment jest podzielony na działy odpowiadające działom realizowanym w klasach szkoły średniej.
  \item Po wejściu do działu użytkownik otrzymuje zadania z całego działu oraz dodatkowe zadania z arkuszy maturalnych.
  \item Zadania otwarte i zamknięte mogą być mieszane w jednej puli.
  \item W zadaniach otwartych użytkownik wpisuje pojedynczą odpowiedź w polu tekstowym; odpowiedź błędna powoduje powrót zadania do puli.
  \item Moduł maturalny nie zawiera odrębnych lekcji. Materiał z rozwiązaniem pojawia się po błędnej odpowiedzi.
  \item Rozwiązania istniejących zadań mogą być współdzielone z materiałami lekcyjnymi, a nowe rozwiązania są potrzebne przede wszystkim dla zadań z arkuszy.
  \item Platforma nie obiecuje "pewniaków maturalnych"; materiał ma systematycznie obejmować pełną podstawę programową.
  \item System treści musi umożliwiać regularne dodawanie autorskich zadań.
\end{itemize}
\end{ElearnRequirement}
\clearpage
\chapter{Skonsolidowany backlog funkcjonalny}

\begin{longtable}{@{}p{0.08\textwidth}p{0.27\textwidth}p{0.45\textwidth}p{0.12\textwidth}@{}}
\toprule
\textbf{ID} & \textbf{Obszar} & \textbf{Zakres} & \textbf{Priorytet}\\
\midrule
\endhead
F-01 & Rejestracja & Konto, potwierdzenie e-mail, powrót do zakupu po aktywacji. & MVP\\
F-02 & Plan nauki & Dni, godzina, dzienny cel liczony zadaniami, zgody na przypomnienia. & MVP\\
F-03 & Nauka & Hierarchia przedmiot -> poziom -> klasa -> dział -> zagadnienie. & MVP\\
F-04 & Zadania & Pula zadań, odpowiedzi A--D, zadania otwarte, mechanizm ponownych prób. & MVP\\
F-05 & Materiały & Lekcje z głosem, rozwiązania z narracją, notatki PDF. & MVP\\
F-06 & Postępy & Liczba poprawnych odpowiedzi względem wszystkich zadań w dziale. & MVP\\
F-07 & Profil & Konto, hasło, kurs i ważność, kody rabatowe, zgody. & MVP\\
F-08 & Kontakt & Formularz e-mail zamiast czatu czasu rzeczywistego. & MVP\\
F-09 & Rabaty & Reguły 5/10/15/20\%, przechowywanie kuponów w profilu. & MVP\\
F-10 & Matura & Zadania działowe i arkuszowe, otwarte i zamknięte, bez osobnych lekcji. & V1\\
F-11 & Opinie & Ankieta dotycząca ulepszeń i repozytorium zgłoszeń. & V1\\
F-12 & Edytor & Diagramy, geometria, wykresy funkcji, eksport zasobów. & V1/V2\\
F-13 & Aplikacja mobilna & Odnośniki do sklepów i dostosowanie odtwarzacza. & V2\\
\bottomrule
\end{longtable}

\section{Kwestie wymagające decyzji}
\begin{itemize}
  \item Dokładny zakres treści dostępnych bez konta.
  \item Kanały powiadomień: wyłącznie e-mail czy także powiadomienia przeglądarkowe i mobilne.
  \item Format i sposób hostowania nagrań lekcji oraz rozwiązań zadań.
  \item Szczegółowy regulamin rabatów, terminy ważności i zasady odnowienia abonamentu.
  \item Miejsce przechowywania odpowiedzi z ankiety "Zagłosuj na ulepszenie".
  \item Wybór edytora rysunków i wykresów oraz wymagania dotyczące eksportu.
  \item Sposób oceniania równoważnych odpowiedzi w zadaniach otwartych, np. ułamków, jednostek i różnych formatów zapisu.
\end{itemize}
\begin{ElearnImportant}
Niniejszy dokument jest punktem wyjścia do przygotowania makiet docelowych, modeli danych, reguł biznesowych i kryteriów akceptacji. Przed implementacją każda pozycja backlogu powinna otrzymać właściciela, priorytet oraz testowalne warunki odbioru.
\end{ElearnImportant}

\appendix
\chapter{Mapa obrazów i komentarzy}
\begin{longtable}{@{}p{0.12\textwidth}p{0.48\textwidth}p{0.28\textwidth}@{}}
\toprule
\textbf{Obraz} & \textbf{Zakres} & \textbf{Komentarz źródłowy}\\
\midrule
\endhead
01 & Dni przypomnień i rejestracja & 0\\
02--05 & Cel dzienny, zgody i podsumowanie & 1--2\\
06 & Panel główny i rabaty & 3\\
07 & Opinie, kontakt, zapisane materiały & 4\\
08 & Struktura sekcji Nauka & 5\\
09 & Widok zagadnienia i zadania & 6\\
10 & Notatki PDF & 7\\
11 & Wyjaśnienia rozwiązań & 8\\
12 & Odpowiedzi i ponowne próby & 9\\
13 & Postępy & 10\\
14 & Lista lekcji & 11\\
15 & Nagranie lekcji & 12\\
16 & Profil użytkownika & 13\\
17--18 & Diagramy i wykresy & 14--15\\
19 & Moduł maturalny & 16\\
\bottomrule
\end{longtable}

\end{document}
